\section{Introduction}
\label{sec:introduction}

A chest-radiograph classifier that knows when to defer is more useful in a
hospital than one that is merely accurate. Selective prediction formalises
this: the model returns a decision only when its confidence clears a threshold,
and refers the remainder to a human reader. The threshold fixes an operating
point --- a coverage --- and with it the workload the system imposes on
radiology.

That operating point is chosen once, on data from the hospital where the system
was developed. It is then deployed somewhere else. Whether it survives the move
is an empirical question that is rarely asked, and the way it is usually asked
answers something narrower than intended: models are re-evaluated at a second
site, but the abstention threshold is re-fitted there too. A re-fitted threshold
tests whether a \emph{method} transfers. It cannot test whether a \emph{deployed
policy} transfers, because the policy that would actually ship is the frozen
one.

A second shift travels with the first and is usually left implicit. Multimodal
chest-radiograph models increasingly consume clinical context alongside the
image --- the indication, the referring history, the reason the study was
ordered. This text is not a stable input. Its availability is a property of how
a particular hospital documents care, and it varies across institutions in ways
that have nothing to do with the patient. A model that learned to lean on the
indication field at one site may find it thin, absent, or differently written at
the next.

Institution and context therefore shift together. Studies that vary one hold the
other fixed: cross-institution work is typically image-only, and multimodal
context work is typically single-institution. Holding one fixed cannot reveal
whether they interact, and interaction is precisely what determines deployment
risk. If a multimodal model is more accurate at the source site \emph{because}
of context, and context is scarcer at the target site, then the measured benefit
and the unmeasured fragility are the same phenomenon observed from two sides.

This paper contributes an evaluation protocol that crosses the two shifts
deliberately, so their interaction can be estimated rather than assumed away:

\begin{itemize}
  \item \textbf{Two sites.} A source hospital, where models are developed and
    the selective policy is calibrated, and an external hospital used only for
    evaluation.
  \item \textbf{Three controlled context states at each site.} Context left
    intact, context removed, and context replaced with another patient's ---
    the last separating \emph{absence} of information from \emph{wrong}
    information.
  \item \textbf{A frozen policy.} Per-pathology thresholds and abstention
    cutoffs are selected once at the source site and transferred unchanged, so
    what is evaluated is the policy that would actually ship.
  \item \textbf{Preregistration} in a machine-readable, hash-manifested
    protocol, with every amendment dated and its confirmatory status recorded,
    so the boundary between prespecified and post-hoc is auditable rather than
    asserted.
\end{itemize}

The contribution is the protocol and what it measures. We propose no new
architecture and no new uncertainty algorithm; the backbones are deliberately
conventional, held constant so that architecture cannot serve as a competing
explanation for any effect observed. This is not a prospective clinical study or
a regulatory validation.

Two further contributions emerged from building the protocol and are reported
because they generalise beyond it. First, a structural hazard in deriving
pre-diagnostic context from radiology reports: a substantial minority of reports
embed a preliminary interpretation inside the report body, so any pipeline that
takes ``everything before the findings section'' as context silently ingests
post-diagnostic text (\cref{sec:report-structure}). Second, a property of the
decision rule this literature commonly uses --- requiring both that a confidence
interval exclude zero and that the point estimate reach a materiality threshold
--- which caps statistical power at \SI{50}{\percent} when the true effect
equals that threshold, regardless of sample size (\cref{sec:power}).
